% =============================================================================
% Section 6: White Holes and Black Holes of Cosmological Origin
% Merged from: agent_26.tex (S6.1), agent_27.tex (S6.2--6.4)
% =============================================================================

\section{White Holes and Black Holes of Cosmological Origin}
\label{sec:6}

%%%%%%%%%%%%%%%%%%%%%%%%%%%%%%%%%%%%%%%%%%%%%%%%%%%%%%%%%%%%%%%%%%%
%% §6.1 White Holes, Grey Holes, and Black Holes
%%%%%%%%%%%%%%%%%%%%%%%%%%%%%%%%%%%%%%%%%%%%%%%%%%%%%%%%%%%%%%%%%%%

\subsection{White Holes, Grey Holes, and Black Holes}
\label{sec:6.1}

All black holes studied in previous sections were regarded as remnants of the
relativistic collapse of massive stars. However, the formation of stars from rarefied
gas is possible only in relatively late stages of the evolution of the Universe
($t \gtrsim 10^7$~years). In earlier stages, according to current cosmological models,
radiation pressure impeded the growth of condensations; so individual bodies
could not form from \emph{small} perturbations of the primeval gas. In the absence of
star formation, there correspondingly should have been no black-hole formation.
However, this conclusion might be wrong. It is quite possible that near the
beginning of the cosmological expansion the matter distribution and metric were
\emph{highly} inhomogeneous and anisotropic. In this case the condensation of individual
bodies would have been possible. Moreover, because in the early stages of the hot
universe the mass-energy in radiation and in particle-antiparticle pairs greatly
exceeded that in baryons, bodies which condensed then must have been made
primarily of photons and pairs. However, because such an ultrarelativistic gas has
an adiabatic index $\Gamma \leq 4/3$, any body made from it is unstable against gravita\-
tional collapse. Thus, one is led to imagine primordial condensations that, like
Wheeler's~(1955) \cite{Wheeler1955} geons, were made primarily from photons and pairs, but unlike
geons, quickly collapsed to form black holes.

Another possible type of strong inhomogeneity in the early universe is a
``white hole'' (Novikov 1964; Ne'eman 1965) \cite{Novikov1964, Neeman1965}. Suppose that some isolated
portions of the universe failed to begin their expansion at the same moment as
the rest of the universe. As measured by external observers, the delay in the
initial expansion of a given region (or ``core'') might be arbitrarily great; and the
delays might vary from core to core. When a ``lagging core'' eventually begins to
expand, and emerges through its gravitational radius, external observers should
see an explosion with a release of tremendous energy. Such an object is, in some
sense, a concrete realization of Ambartsumyan's (1961, 1964) concept of a super\-
dense ``D-body''. (Thus, one can regard ``D-body'', ``white hole'', and ``lagging
core'' as different names for the same concept.)

Finally, if a lagging core, when it begins to expand, has insufficient energy to
emerge through its gravitational radius into the external universe, then one can
regard it as a ``grey hole''. Thus, in the idealized spherical case, the matter of a
grey hole is ``born'' in the Kruskal diagram, staying always to the left of the center line; it moves
upward in the Kruskal diagram, staying always by a past horizon or a future
horizon; and it finally dives to its death in the terminal $r = 0$ singularity. What
would all three types of cosmological holes actually exist in Nature? What
existence be their properties? Theorists have very little observational way to rule out their
about them.

In these notes we shall dwell on only two points: (i) the growth of cosmological
holes by accretion, and (ii) a limit on how many baryons can be inside cosmological
holes.

%% ----------------------------------------------------------------
%% 6.2  The Growth of Cosmological Holes by Accretion
%% ----------------------------------------------------------------

\subsection{The Growth of Cosmological Holes by Accretion}
\label{sec:6.2}

If cosmological holes have existed since near the beginning of the universe,
then during the era when the energy density in radiation exceeded that in matter,
accretion of radiation onto the holes must have been important (Zel'dovich and
Novikov 1966).  For a rough estimate of the effects of such accretion we can use
equation~(4.2.14) for the accretion of gas onto a stationary hole, taking the
velocity of sound to be $a_\infty = c/\sqrt{3}\,c$:
%
\begin{equation}
\tag{6.2.1}
dM/dt \simeq r_g^2 \rho^{\text{rad}}\!.
\end{equation}
%
In standard cosmological models the radiation energy density varies as
$\rho^{\text{rad}} = 1/Gt^2$.

Putting this into the growth equation~(6.2.1), and integrating from the moment
$t_0$ when a hole of initial mass $M_0$ first forms, we obtain for the final mass, at
$t > t_0$:
%
\begin{equation}
\tag{6.2.2}
M = M_0 (1 - G M_0 / c^3 t_0)^{-1}.
\end{equation}

Notice that the final mass diverges for $t_0 \to GM_0/c^3$.  However, the method of
calculation breaks down for this same limit because, for $t_0 \sim GM_0/c^3$ the
characteristic time scale for the accretion is the same as that for the change of
$\rho^{\text{rad}}$, so the ``steady-state'' hypothesis on which equation
(6.2.1) is based.  To determine whether the accretion is catastrophically great at
early times ($t_0 \lesssim GM_0/c^3$), one must solve the accretion problem in a nonsteady,
cosmological context.  One would not be surprised if the final answer depends
sensitively on the particular choice of initial conditions ($t_0 = GM_0/c^3$ or
$t_0 = 0.01\, GM_0/c^3$ or $t_0 = 10^{-10} GM_0/c^3$).  But the problem has not yet been
solved.

%% ----------------------------------------------------------------
%% 6.3  Limit on the Number of Baryons in Cosmological Holes
%% ----------------------------------------------------------------

\subsection{Limit on the Number of Baryons in Cosmological Holes}
\label{sec:6.3}

Independently of the issue of accretion, one can place a tight limit on what
fraction~$\alpha$ of all baryons in the Universe were swallowed into cosmological holes
in the early stages.

Suppose that prior to some particular early moment~$t_*$, a fraction~$\alpha$ of all
baryons had been swallowed into holes.  At the moment~$t_*$, the ratio of rest mass-energy in baryons to mass-energy in radiation and pairs was tiny
%
\begin{equation}
\tag{6.3.1}
\beta \equiv \left[\frac{\rho_0}{\rho^{\text{rad}} + \rho^{\text{(pairs)}}}\right]_{t_*}
\sim 10^{-7} (t_*/1\;\text{sec})^{1/2} \ll 1.
\end{equation}
%
However, as the universe subsequently expanded, this ratio increased until today
it is far greater than unity.  Moreover, the total mass-energy in holes divided by
the volume of the universe ($\alpha \rho_0$) changed in the same manner as
$\rho_0$ ($\propto$~vol$^{-1}$), if no new holes were formed after~$t_*$, and if accretion was negligible.
Otherwise it increased relative to~$\rho_0$.  This means that
%
\begin{equation}
\tag{6.3.2}
\frac{\alpha}{\beta(1-\alpha)}
\leq \left[\frac{\rho^{\text{(holes)}}}{\rho_0}\right]_{t_*}
\leq \left[\frac{\rho^{\text{(holes)}}}{\rho_0}\right]_{\text{today}}
< 80
\end{equation}
%
Here $\alpha/[\beta(1-\alpha)]$ is the value of $[\rho^{\text{(holes)}}/\rho_0]_{t_*}$ if the holes were all created from
primeval plasma at the moment~$t_*$; if they were created even earlier, when rest
mass was even less important, $[\rho^{\text{(holes)}}/\rho_0]_*$ would be even bigger.  The number
80 is a generous observational upper limit on the ratio of nonluminous matter
to luminous matter in the universe today.

Equation~(6.3.2) for $\beta \ll 1/80$ can be rewritten as
%
\begin{equation}
\tag{6.3.3}
\alpha < 1/80\beta \sim 10^{-5} (t_*/1\;\text{sec})^{-1/2}\,.
\end{equation}

This is a very tight limit on the amount of rest mass that could have been down
black holes at early times~$t_*$.

One can also place a tight limit on the amount of mass-energy that white holes
have spewed forth as radiation in recent times (at cosmological redshifts $z \lesssim 10$).
That mass-energy cannot exceed the total mass-energy in radiation today,
%
\begin{equation}
\tag{6.3.4}
\rho_{\text{from white holes}} < \rho_0^{\text{(rad)}}
= 5 \times 10^{-13}\;\text{erg/cm}^3
= 6 \times 10^{-34}\;\text{g/cm}^3\,.
\end{equation}

%% ----------------------------------------------------------------
%% 6.4  Caution
%% ----------------------------------------------------------------

\subsection{Caution}
\label{sec:6.4}

We conclude with a word of caution.  This cosmological section has dealt with
objects about which both theory and observation are equivocal.  There is no firm
theoretical reason for believing that cosmological holes exist, though theory
surely permits them.

By contrast, theory virtually demands the existence of black holes formed
by stellar collapse.  It would be astonishing indeed if every star in the Galaxy
somehow managed to avoid the black-hole fate!  [See Hoyle, Fowler, Burbidge,
and Burbidge (1964).]

The clear discovery of black holes, we expect, is only a few months or years
away.  (The X-ray source Cyg~X-1 is an excellent candidate.)  But cosmological
holes might never be found and might, in fact, not exist at all.  On the other
hand, we might eventually discover that they are the energizers of quasars and
of explosions in galactic nuclei.
